\documentclass[11pt]{article}

\newcommand{\cnum}{Math 245B}
\newcommand{\ced}{Winter 2026}
\newcommand{\ctitle}[4]{\title{\vspace{-0.5in}\cnum, \ced\\Assignment #1: #2}\author{\vspace{-0.35in}\\#3, #4}}
\usepackage{enumitem}
\usepackage[usenames,dvipsnames,svgnames,table,hyperref]{xcolor}
\usepackage{amsmath, amsfonts}
\usepackage{graphicx} % Required for inserting images
\usepackage{float}
\usepackage{listings}
\usepackage{xcolor}
\usepackage{hyperref}
\usepackage{comment}

\renewcommand*{\theenumi}{\alph{enumi}}
\renewcommand*\labelenumi{(\theenumi)}
\renewcommand*{\theenumii}{\roman{enumii}}
\renewcommand*\labelenumii{\theenumii.}

\definecolor{codegreen}{rgb}{0,0.6,0}
\definecolor{codegray}{rgb}{0.5,0.5,0.5}
\definecolor{codepurple}{rgb}{0.58,0,0.82}
\definecolor{backcolour}{rgb}{0.95,0.95,0.92}
\definecolor{header}{HTML}{205459}
\definecolor{solution}{HTML}{204d52}

\newcommand{\solution}[1]{{{\textcolor{header}{Solution:} \textcolor{solution}{#1}}}}
\setlist[enumerate]{label=(\alph*)}

\lstdefinestyle{mystyle}{
    backgroundcolor=\color{backcolour},
    commentstyle=\color{codegreen},
    keywordstyle=\color{magenta},
    numberstyle=\tiny\color{codegray},
    stringstyle=\color{codepurple},
    basicstyle=\ttfamily\footnotesize,
    breakatwhitespace=false,
    breaklines=true,
    captionpos=b,
    keepspaces=true,
    numbers=left,
    numbersep=5pt,
    showspaces=false,
    showstringspaces=false,
    showtabs=false,
    tabsize=2
}


\begin{document}
\ctitle{\#1}{}{Asher Christian}{006-150-286}
\date{}
\maketitle
\section{Kolomogorov 0-1 Law}
\emph{
    For independent $\{X_i\}_{i=1}^{\infty}$ if $A \in \mathcal{T}$ then  $ \mathbb{P}(A) = 0$ or  $ \mathbb{P}(A) = 1$\\
}
(try to prove $A \perp A$.  $ \mathbb{P}(A) = 1, \text{or}, \mathbb{P}(A) = 0 \implies A \perp \mathcal{F}$. If $A \perp A \implies \mathbb{P}(A) = \mathbb{P}(A)^2$.
$ \mathcal{T}$ the tail sigma algebra
 \[
     \bigcap_{n=1}^{\infty}\sigma(X_n,X_{n+1},\dots)
\] 
\[
    \mathcal{T} \perp \sigma(X_1,X_2) \qquad \mathcal{T} \perp \mathcal{F}_n = \sigma(X_1,\dots,X_n)
\] 
want to show
\[
    \mathcal{F} = \sigma(X_1,\dots,X_n,\dots) = \sigma(\bigcup_{n=1}^{\infty} \mathcal{F}_n)
\] 
\[
    \mathcal{T} \perp \mathcal{T} \impliedby \mathcal{T} \perp \mathcal{F} \impliedby? \;\; \mathcal{T} \perp \bigcup_{n=1}^{\infty} \mathcal{F}_n
\] 
We know $\mathcal{T} \perp (\bigcup_{n=1}^{\infty} \mathcal{F}_n$ does this imply $\mathcal{T} \perp \sigma(\bigcup_{n=1}^{\infty}\mathcal{F}_n)$ $\pi-\lambda$ again
\begin{enumerate}
    \item show that $\bigcup_{n=1}^{\infty} \mathcal{F}_n$ is a $ \pi $-system
    \item $\{B : B \perp \mathcal{T}\}$ is a  $\lambda$-system
    \item $\pi-\lambda \implies \sigma(\bigcup_{n=1}^{\infty} \mathcal{F}_n \perp \mathcal{T}$
\end{enumerate}

\section{Percolation}
Infinite lattice $ \mathbb{Z}^2$
$0 \le p \le 1$  $X_e \sim \text{Bern}(p)$ if it is zero the edge is destroyed if it is 1 the edge is maintained
$ \mathbb{P}( C_{\infty} \text{exists})$ the probability that there is a connected component that it is infinitely large $|CP| = +\infty$.
This is a function of $p$ Let  $A = \exists CP_{\infty}$ then $A \perp X_{e_1}, X_{e_2}$ $A \in \mathcal{T}$ so  $f(p) = 0$ or $f(p) = 1$ The jump point is
when  $d = 2, p_c = \frac{1}{2}$ 

\section{H-S 0-1 Law}
Hewitt Savage
\[
    A := \{\lim_{n\to \infty}\frac{S_n}{\sqrt{n}} \le a\}, \qquad A \in \mathcal{T}
\] 
\[
S_n = \sum_{i=1}^{n}X_i
\] 
\[
    \tilde{A} := \{ \lim_{}S_n \le a\} \qquad \tilde A \not\in \mathcal{T}
\] 
\[
\lim_{n\to \infty}S_n(X_1,\dots,X_n,\dots) = \lim_{n\to \infty} S_n(X_2,X_1,X_3,\dots,X_n,\dots)
\] 
This limit function is symmetric for example
\[
\sum_{i=1}^{\infty}X_i \qquad \sum_{i,j}^{}X_iX_j
\] 
\[
\vec X := (X_1,\dots,X_n,\dots)
\] 
\[
    \vec X : (\Omega, \mathcal{F}) \rightarrow ( \mathbb{R}^{\infty}, \mathcal{R}^{\infty})
\] 
\emph{ Definition:\\
$A \in \mathcal{R}^{\infty}$  is symmetric $\iff$  $\forall i < j$ $(a_1,a_2,\dots,a_n,\dots) \in A \implies (a_1,\dots,a_j, \dots, a_i, \dots, a_n, \dots) \in A$
}\\
What about antisymmetric like determminant\\
\emph{
    H-S 0-1 Law:\\
    If $\{X_i\}_{i=1}^{\infty}$ are i.i.d R.V.  $A \in \mathcal{R}^{\infty}$ is symmetric then $  \mathbb{P}(\vec X \in A) = 0$ or $1$
}
Example
\[
    \mathbb{P}(\lim_{n\to \infty}S_n \le a) = 0 \;\; \text{or} \;\; 1
\] 
.
If $X_i$ are i.i.d r.v. if  $E \in \mathcal{E} \subset \mathcal{F}^{\infty}$ $E$ is exchangable then  $ \mathbb{P}(\vec X \in E) = 0$ or $1$
\section{Lemma}
\emph{
$X_i$ i.i.d,  $ \mathbb{P}(X_1 = 0) \ne 1$ then 
\[
    \limsup_{n\to \infty} S_n = -\infty \;\; \text{or} \;\; +\infty
\] 
}\\
Proof: $\forall$ fixed  $L \in \mathbb{R}$
\[
    \mathbb{P}(\limsup_{n\to \infty}S_n \le L) = 0 \;\; \text{or} \;\; 1 = g(L)
\] 
If  $g(L)$ is sconstant 0 or 1 then we are done.  otherwise $\exists L_c$ such that  $g(L) = 1$ when $ L < c$  $g(L) = 0$ when  $ L > L_c$ then
 \[
     \limsup_{n\to \infty} S_n = L_c \qquad \text{a.s.}
\] 
Assume $\exists L_c$ such that  $\limsup_{ n\to \infty } S_n = L_c$ almost sure. $\limsup_{n\to \infty}S_n  = X_1 + \limsup_{n\to \infty}\tilde S_n$ 
\[
    (\tilde X_1, \tilde X_2, \dots) = (X_2, X_2,\dots)
\] 
\[
    \tilde {\vec X} \sim^{d} \vec X
\] 
so
\[
\limsup_{n\to \infty} S_n = X_1 + \limsup_{n\to \infty} S_n
\] 
so
\[
    L_c = X_1 + L_c  \qquad \text{a.s.}
\] 
so
\[
    X_1 = 0 \qquad \text{a.s}
\] 
contradiction.
\section{Proof of H-S Law}
\[
    E \subset \mathcal{E} \subset \mathcal{F}^{\infty}
\] 
then for any $\epsilon > 0$ there exists $n, A \in \mathcal{F}_n$ such that 
$ \mathbb{P}(E \triangle A) < \epsilon$ 
\emph{
    Lemma: $ \mathcal{A}$ is an algebra, $ \mathcal{F} = \sigma( \mathcal{A})$. 
    \[
    \forall B \in \mathcal{F} \;\;\; \delta > 0 \;\; \exists \; A \in \mathcal{A} \rightarrow \mathbb{P}(A \triangle B) < \delta
    \] 
    Proof using $\pi-\lambda$ 
    $ \mathcal{P} := \mathcal{A}$
    $ \mathcal{L} := \{ B \in \mathcal{F}: \forall \delta , \exists A \in \mathcal{A}, \mathbb{P}(A \triangle B) < \delta\}$
    Then $\mathcal{P} \subset \mathcal{L}$ 
    \begin{enumerate}
        \item $B \in \mathcal{L}, B^{c} \in \mathcal{L}$ since $ \mathcal{A}$ is closed under complement
        \item $\bigoplus_{i=1}^{\infty} B_i \in \mathcal{L}$ Fix $\delta$ cut off the tail so that the head is all the probability except for $\frac{\delta}{3}$ and find $A \in \mathcal{A}$ that approximate the first $B_i$ then the union of these  $A_i$ is close to the first  $B_i$
    \end{enumerate}
}
$A \in \mathcal{F}_n$ $ A \in \sigma(X_1,X_2,\dots,X_n)$
$\tilde A \in \sigma(X_{n+1},\dots X_n)$ $\tilde A = \Pi \circ A$
 \[
\Pi = 
\begin{cases}
    X_1 \rightarrow X_{n+1}\\
    X_2 \rightarrow X_{n+2}\\
    \dots\\
    X_n \rightarrow X_{2n}
\end{cases}
\] 
Then
 \[
\mathbb{P}(\vec X \in E \triangle \tilde A) = \mathbb{P}(\vec X \in E \triangle A)
\] 


\end{document}
